% ══════════════════════════════════════════════════════════════════════════════
%  EXPERIMENT 2 — A COOLING CURVE                                       (Lesson 6)
%  Content from the A4 master's Lesson 9. First line graph of the year.
% ══════════════════════════════════════════════════════════════════════════════
\experimentdiv{A Cooling Curve}{Temperature over time \textbullet{} Line graphs \textbullet{} The freezing point \textbullet{} The plateau}{6}\label{sec:2}
\lessoncodes{VC2S8U05 \textbullet{} VC2S8I01--I06}

\vocabhead
\begin{ctab}
\begin{tabular}{|SW{4.2cm}|W{11.6cm}|}
  \hline\hd & \thead{Word} & \thead{Definition (copy from the slide)} \\ \hline
  & \scaf{\textbf{Cooling curve}} & \scaf{A line graph showing how the temperature of a substance falls over time.} \\ \hline
  & \textbf{Freezing point} & \\ \hline
  & \textbf{Line graph} & \\ \hline
  & \textbf{Axis} & \\ \hline
  & \textbf{Plateau} & \\ \hline
\end{tabular}
\end{ctab}

\begin{lipanel}
\inquiryli{Record temperature every minute, then draw a line graph and describe the pattern.}
\understandli{Find the freezing point from the graph, and explain why the temperature stays constant while a substance freezes.}
\sctitle
\begin{itemize}[leftmargin=1.4em, itemsep=0pt, topsep=1pt]
  \item I can record the temperature every minute and draw a line graph with labelled axes, even scales and a title.
  \item I can read the freezing point from my graph and explain the plateau using particles.
\end{itemize}
\end{lipanel}

\begin{skillsfocus}
\focus{QP1}{Sketch how the temperature will change.}
\focus{PC3}{Assess the risks before you start.}
\focus{PC6}{Read the temperature every minute.}
\focus{PA2}{Draw your first line graph.}
\focus{PA5}{Describe the curve's shape.}
\focus{PA8}{Explain the plateau with particles.}
\focus{EV2}{Compare with the accepted 44\,\textdegree C.}
\end{skillsfocus}

% ─── TECHNIQUE ────────────────────────────────────────────────────────────────
\exppart{Technique spotlight}{Blue}

\begin{technique}{Drawing a line graph}{PA2}
\begin{minipage}[c]{0.56\linewidth}
\begin{enumerate}[leftmargin=1.6em, itemsep=1pt, topsep=0pt]
  \item Put the thing you \textbf{changed or controlled} (time) on the horizontal axis.
  \item Put the thing you \textbf{measured} (temperature) on the vertical axis.
  \item Choose an \textbf{even scale} for each axis that uses most of the grid.
  \item \textbf{Label} both axes with the quantity and the unit, e.g.\ \emph{time (minutes)}.
  \item \textbf{Plot} each point accurately with a small cross, and give the graph a \textbf{title} that names both variables.
\end{enumerate}
\end{minipage}\hfill
\begin{minipage}[c]{0.4\linewidth}\centering
\begin{tikzpicture}[x=0.5cm, y=0.05cm, font=\scriptsize]
  \foreach \m in {0,...,8}{\draw[black!25] (\m,20) -- (\m,70);}
  \foreach \t in {20,30,...,70}{\draw[black!25] (0,\t) -- (8,\t); \node[anchor=east] at (0,\t) {\t};}
  \foreach \m in {0,2,...,8}{\node[anchor=north] at (\m,20) {\m};}
  \draw[axis] (0,20) -- (8.5,20); \draw[axis] (0,20) -- (0,74);
  \draw[Ink, line width=1.2pt] plot[smooth] coordinates {(0,66) (1,55) (2,48) (3,44) (4,44) (5,44) (6,44) (7,40) (8,36)};
  \foreach \p in {(0,66),(1,55),(2,48),(3,44),(4,44),(5,44),(6,44),(7,40),(8,36)} \node[font=\small] at \p {$\times$};
  \node[anchor=north] at (4,13) {time (minutes)};
  \node[rotate=90, anchor=south] at (-1.1,45) {temperature (\textdegree C)};
\end{tikzpicture}
\end{minipage}
\end{technique}

\begin{technique}{Join the points, or a line of best fit?}{PA2}
\begin{minipage}[c]{0.56\linewidth}
\textbf{Joining the points} (left): a ruled straight line from each point to the next. It shows every reading exactly, but each small measuring error becomes a kink, and it suggests you know what happened between readings when you do not.

\smallskip
\textbf{A line of best fit} (right): \emph{one} smooth line or curve through the middle of the points, with about as many points above it as below. It ignores an anomaly, evens out small errors and shows the \textbf{trend}. Scientists almost always use it.

\smallskip
Start by joining the points if you need to, but from your cooling curve on, \textbf{draw a line of best fit}.
\end{minipage}\hfill
\begin{minipage}[c]{0.4\linewidth}\centering
\begin{tikzpicture}[x=0.3cm, y=0.3cm, font=\scriptsize]
  \foreach \ox/\lab in {0/joined, 7.5/best fit}{
    \begin{scope}[shift={(\ox,0)}]
    \draw[axis] (0,0) -- (6.2,0); \draw[axis] (0,0) -- (0,6.2);
    \node[anchor=north] at (3,-0.2) {\lab};
    \end{scope}}
  \draw[Ink, line width=0.8pt] (0.5,0.9) -- (1.5,2.1) -- (2.5,2.4) -- (3.5,3.9) -- (4.5,4.1) -- (5.5,5.4);
  \draw[Ink, line width=1.1pt] (8.0,0.8) -- (13.0,5.3);
  \foreach \x/\y in {0.5/0.9,1.5/2.1,2.5/2.4,3.5/3.9,4.5/4.1,5.5/5.4}{
    \node[font=\small] at (\x,\y) {$\times$};
    \node[font=\small] at (\x+7.5,\y) {$\times$};}
\end{tikzpicture}
\end{minipage}
\end{technique}

\begin{workedex}{Choosing a scale}
Temperatures in a set of results go from 22\,\textdegree C up to 68\,\textdegree C. The grid has 15 big squares going up. Choose a scale for the temperature axis.

\medskip
Start just below the smallest value, at \textbf{20\,\textdegree C}. Going up in steps of \textbf{5\,\textdegree C} per big square, 15 squares reach $20 + 15 \times 5 = 95$\,\textdegree C. That fits every result and uses most of the grid.
\end{workedex}

\begin{wedo}{Choosing a scale}
Times in a set of results go from 0 to 12 minutes. The grid has 15 big squares going across. Choose a scale for the time axis, and show it fits.
\workbox{3cm}
\end{wedo}

\begin{youdobox}[Your first line graph]
\task A student heated 150\,mL of water and recorded its temperature every minute. Draw a line graph of the results. Label both axes with units, give the graph a title, then join the points with a ruler. Would a line of best fit look different here? \skr{PA2}

\begin{ctab}
\begin{tabular}{|l|*{11}{C{0.78cm}|}}
  \hline\thead{Time (min)} & 0 & 1 & 2 & 3 & 4 & 5 & 6 & 7 & 8 & 9 & 10 \\ \hline
  \thead{Temperature (\textdegree C)} & 18 & 27 & 36 & 45 & 53 & 61 & 69 & 77 & 84 & 91 & 97 \\ \hline
\end{tabular}
\end{ctab}
\graphgrid{15cm}
\end{youdobox}

% ─── PLAN AND SAFETY ─────────────────────────────────────────────────────────
\exppart{Plan and safety}{Teal}

\textbf{Aim:} To measure the temperature of lauric acid as it cools and freezes, and find its freezing point.

\begin{infobox}{Why lauric acid?}
Water would need a freezer to freeze. Lauric acid, a white waxy solid found in coconut oil, melts at a temperature hot water from a kettle can reach, so you can melt it and then watch it cool and freeze on the bench.
\end{infobox}

\needspace{14\baselineskip}
\task Sketch what you think the temperature of the lauric acid will do as it cools, then explain your sketch. Will it fall steadily the whole time? \skr{QP1}

\noindent\begin{minipage}[t]{0.42\linewidth}\drawbox{5cm}{\scriptsize temperature $\uparrow$ \hfill time $\rightarrow$}\end{minipage}\hfill
\begin{minipage}[t]{0.55\linewidth}\answerbox{5cm}\end{minipage}

\needspace{14\baselineskip}
\task Complete the risk assessment before you start. For each risk, say how to prevent it and what to do if it happens. Your teacher must initial it. \skr{PC3}\par\nobreak\hfill\textbf{Teacher's initials:}\ \underline{\hspace{3cm}}
\begin{ctab}
\begin{tabular}{|SW{3.8cm}|W{5.4cm}|W{5.8cm}|}
  \hline\hd & \thead{Risk (what could go wrong)} & \thead{Prevention} & \thead{Responding (if it happens)} \\ \hline
  & \scaf{Scalds from the hot-water bath} & \scaf{Lift the boiling tube with a test-tube holder, never with fingers} & \scaf{Cool the burn under running water for 20 minutes; tell the teacher} \\ \hline
  & Lauric acid gets on skin or in eyes & & \\ \hline
  & The glass thermometer breaks & & \\ \hline
  & Burns from the hot boiling tube & & \\ \hline
\end{tabular}
\end{ctab}

% ─── METHOD ──────────────────────────────────────────────────────────────────
\exppart{Method}{Teal}

\subsubsection*{Equipment}
\begin{multicols}{2}
\begin{itemize}[leftmargin=1.2em, itemsep=0pt, topsep=0pt]
  \item boiling tube of lauric acid, melted by your teacher in a hot-water bath
  \item test-tube holder
  \item thermometer ($-10$ to 110\,\textdegree C)
  \item retort stand and clamp
  \item 250\,mL beaker (empty)
  \item stopwatch
  \item safety glasses
\end{itemize}
\end{multicols}

\begin{center}
\begin{tikzpicture}[scale=0.95]
  \begin{scope}\eqStand\end{scope}
  \draw[eq] (2.2,0.2) -- (2.2,0.21);
  \draw[eq] (1.55,2.25) -- (1.55,0.6) arc[start angle=180,end angle=360,radius=0.3] -- (2.15,2.25);
  \fill[black!12] (1.57,1.5) -- (1.57,0.6) arc[start angle=180,end angle=360,radius=0.28] -- (2.13,1.5) -- cycle;
  \draw[Ink, line width=0.6pt] (1.57,1.5) -- (2.13,1.5);
  \draw[eq] (1.05,2.05) -- (1.25,1.9) -- (1.25,0.2) -- (2.45,0.2) -- (2.45,1.95);
  \draw[eq] (1.85,0.65) circle (0.1);
  \draw[eq] (1.79,0.74) -- (1.79,4.3) (1.91,0.74) -- (1.91,4.3);
  \eqlab{(1.91,3.9)}{(3.4,3.9)}{thermometer}
  \eqlab{(1.95,3.0)}{(3.4,3.2)}{clamp}
  \eqlab{(2.15,2.0)}{(3.4,2.4)}{boiling tube}
  \eqlab{(2.13,1.2)}{(3.4,1.6)}{lauric acid}
  \eqlab{(2.45,0.9)}{(3.4,0.8)}{empty beaker}
  \eqlabl{(0.35,3.6)}{(-0.4,3.6)}{retort stand}
\end{tikzpicture}
\expfigcap{Scientific diagram of the cooling curve set-up.}
\end{center}

\begin{enumerate}[itemsep=2pt]
  \item Put on safety glasses. Collect a boiling tube of melted lauric acid from the hot-water bath, using a test-tube holder.
  \item Clamp the boiling tube so it stands upright inside the empty beaker.
  \item Put the thermometer in the lauric acid, with its bulb in the middle of the liquid.
  \item When the temperature reads 65\,\textdegree C or just below, start the stopwatch and record the temperature.
  \item Record the temperature every minute for 12 minutes. Stir gently with the thermometer while the lauric acid is still liquid.
  \item Write down the time when you first see solid forming, and when it has all gone solid.
  \item Do not try to pull the thermometer out once the lauric acid is solid. Return the tube to the hot-water bath to melt it again.
\end{enumerate}

% ─── RESULTS ─────────────────────────────────────────────────────────────────
\exppart{Results}{Amber}

\subsubsection*{The temperature of lauric acid as it cools \skr{PC6}}
\begin{ctab}
\begin{tabular}{|SZ{3.2cm}|Z{3.6cm}|Z{3.2cm}|Z{3.6cm}|}
  \hline\hd & \thead{Time (min)} & \thead{Temperature (\textdegree C)} & \thead{Time (min)} & \thead{Temperature (\textdegree C)} \\ \hline
  & \scaf{e.g. 0} & \scaf{65} & \scaf{e.g. 7} & \scaf{44} \\ \hline
  & 0 & & 7 & \\ \hline
  & 1 & & 8 & \\ \hline
  & 2 & & 9 & \\ \hline
  & 3 & & 10 & \\ \hline
  & 4 & & 11 & \\ \hline
  & 5 & & 12 & \\ \hline
  & 6 & & & \\ \hline
\end{tabular}
\end{ctab}

\noindent Solid first seen at \rule{2.2cm}{0.4pt} minutes. \qquad All solid at \rule{2.2cm}{0.4pt} minutes.

\clearpage
\task Draw a cooling curve of your results. Put \textbf{time (minutes)} on the horizontal axis and \textbf{temperature (\textdegree C)} on the vertical axis. Plot each point, draw a smooth \textbf{line of best fit} (a curve is fine) and give the graph a title. \skr{PA2}

\vspace{3pt}
\graphgrid{16cm}

% ─── ANALYSIS ────────────────────────────────────────────────────────────────
\exppart{Analysis and evaluation}{Coral}

\task Describe the shape of your cooling curve, from start to finish. \skr{PA5}
\answerbox{5cm}

\needspace{14\baselineskip}
\task Use your graph. \skr{EV2}
\pt{a}Read the freezing point of lauric acid from your graph.
\answerbox{1cm}
\pt{b}The accepted freezing point of lauric acid is 44\,\textdegree C. Compare your value with it, and suggest one reason for any difference.
\answerbox{5cm}
\pt{c}For how many minutes did the temperature stay on the plateau?
\workbox{2cm}

\needspace{10\baselineskip}
\task Was your prediction correct? Explain, using your graph. \skr{EV1}
\answerbox{5cm}

\needspace{14\baselineskip}
\task Draw particle diagrams of the lauric acid before, during and after the plateau. Use the same number of particles, the same size, in each. \skr{PA4}

\noindent
\begin{minipage}[t]{0.32\linewidth}\centering\drawbox{4.5cm}{}{\small Before: liquid}\end{minipage}\hfill
\begin{minipage}[t]{0.32\linewidth}\centering\drawbox{4.5cm}{}{\small During: freezing}\end{minipage}\hfill
\begin{minipage}[t]{0.32\linewidth}\centering\drawbox{4.5cm}{}{\small After: solid}\end{minipage}

\needspace{10\baselineskip}
\task Explain, using particles, what was happening to the lauric acid during the plateau, and why the temperature stayed steady. \skr{PA8}
\answerbox{5cm}

\needspace{16\baselineskip}
\task Another group recorded these temperatures every minute, starting at 0 minutes: 64, 57, 51, 54, 44, 44, 44, 44, 43, 40, 37\,\textdegree C. \skr{PA7}
\pt{a}Which reading does not fit the pattern?
\answerbox{1cm}
\pt{b}Suggest what might have caused it.
\answerbox{5cm}
\pt{c}What freezing point do their results show?
\answerbox{1cm}

\needspace{14\baselineskip}
\task Evaluate the method. \skr{EV3}
\pt{a}Identify one source of error in this experiment.
\answerbox{5cm}
\pt{b}Suggest one way to make the results more reliable.
\answerbox{5cm}

\begin{yr8box}
\begin{keepext}
\exttask A group used twice as much lauric acid in their tube. Predict how their cooling curve would be different from yours, and how it would be the same. Explain both.
\answerbox{5cm}
\end{keepext}
\begin{keepext}
\exttask Candle wax is a mixture, not a pure substance. Its cooling curve has no flat plateau, only a part where the line slopes more gently. Suggest why.
\answerbox{5cm}
\end{keepext}
\begin{keepext}
\exttask The lauric acid cooled quickly at the start and slowly near the end. Think about how far its temperature was above room temperature in each case, and suggest a rule for how fast hot things cool.
\answerbox{5cm}
\end{keepext}
\end{yr8box}

% ─── SKILL CHECK ─────────────────────────────────────────────────────────────
\exppart{Skill check}{NavyBlue}

\begin{skillcheck}
\scrow{QP1}{I can sketch and explain a prediction.}
\scrow{PC3}{I completed a risk assessment and followed it.}
\scrow{PC6}{I recorded a reading at every whole minute, with its unit.}
\scrow{PA2}{I can draw a line graph with labelled axes, even scales, plotted points and a title.}
\scrow{PA5}{I can describe the shape of a cooling curve.}
\scrow{PA8}{I can use the particle model to explain the plateau.}
\scrow{EV2}{I can compare my result with the accepted value.}
\end{skillcheck}

\task Which skill will you focus on in Experiment 3? Explain what you will do differently.
\answerbox{5cm}
